\documentclass[11pt]{article}

\usepackage[T1]{fontenc}
\usepackage[utf8]{inputenc}
\usepackage{lmodern}
\usepackage[margin=1in]{geometry}

\usepackage{amsmath}
\usepackage{amssymb}
\usepackage{mathtools}
\usepackage{siunitx}

\usepackage{graphicx}
\usepackage{tikz}
\usetikzlibrary{arrows.meta,positioning}
\usepackage{xcolor}
\usepackage{booktabs}
\usepackage{caption}
\usepackage{subcaption}

\usepackage{listings}

\definecolor{codeKeyword}{RGB}{0,0,140}
\definecolor{codeComment}{RGB}{0,100,0}
\definecolor{codeString}{RGB}{160,20,20}
\definecolor{codeNumber}{RGB}{90,90,90}
\definecolor{codeBackground}{RGB}{248,248,248}

\lstdefinelanguage{Verilog}{
  morekeywords={
    module,endmodule,input,output,inout,wire,reg,logic,assign,always,always_ff,
    always_comb,always_latch,if,else,case,endcase,for,generate,endgenerate,
    begin,end,parameter,localparam,typedef,struct,enum,signed,unsigned
  },
  sensitive=true,
  morecomment=[l]{//},
  morecomment=[s]{/*}{*/},
  morestring=[b]"
}

\lstdefinestyle{thesiscode}{
  backgroundcolor=\color{codeBackground},
  basicstyle=\ttfamily\small,
  numbers=left,
  numberstyle=\scriptsize\color{codeNumber},
  stepnumber=1,
  numbersep=8pt,
  keywordstyle=\color{codeKeyword}\bfseries,
  commentstyle=\color{codeComment}\itshape,
  stringstyle=\color{codeString},
  frame=single,
  framerule=0.4pt,
  breaklines=true,
  breakatwhitespace=true,
  showstringspaces=false,
  tabsize=2,
  keepspaces=true,
  captionpos=b
}

\lstset{style=thesiscode}

\usepackage[hidelinks]{hyperref}
\usepackage[nameinlink,noabbrev]{cleveref}

\newcommand{\SNR}{\ensuremath{\mathrm{SNR}}}
\newcommand{\dBFS}{\ensuremath{\mathrm{dBFS}}}
\newcommand{\ENOB}{\ensuremath{\mathrm{ENOB}}}
\newcommand{\AXIStream}{\ensuremath{\mathrm{AXI4\mbox{-}Stream}}}
\newcommand{\IandQ}{\ensuremath{I\!\!/\!\!Q}}


\title{Monograph 01: Spatial Logic}
\author{Radar PHY and FPGA Architecture Study}
\date{\today}

\begin{document}

\maketitle

\section{Introduction}
This monograph frames the spatial logic perspective of radar gateware realization, from logic-cell granularity up through CLB and die-level organization. The discussion establishes how placement pressure, routing locality, and timing closure shape deterministic throughput for \AXIStream data movement and coherent \IandQ signal transport.

The objective is to map architectural intent into physically realizable structures while preserving signal fidelity metrics such as \SNR, dynamic range in \dBFS, and converter effectiveness in \ENOB.

\section{Temporal vs. Spatial Computing: The Physics of Determinism}

\subsection{The Speed of Light Constraint in Radar Timing}
\textbf{Derivation outline.}
For a monostatic radar, range is inferred from round-trip time-of-flight:
\begin{equation}
R=\frac{ct}{2}.
\end{equation}
A timing perturbation \(\Delta t\) maps directly to a range perturbation:
\begin{equation}
\Delta R=\frac{c\Delta t}{2}.
\end{equation}

\textbf{Numerical checkpoints.}
\begin{itemize}
  \item For \(\Delta t=\SI{100}{\nano\second}\): \(c\Delta t\approx\SI{30}{\meter}\) propagation-path ambiguity, hence \(\Delta R\approx\SI{15}{\meter}\) range ambiguity in monostatic operation.
  \item For \(\Delta t=\SI{1}{\nano\second}\): \(c\Delta t\approx\SI{0.30}{\meter}\), hence \(\Delta R\approx\SI{0.15}{\meter}\).
\end{itemize}
These checkpoints motivate nanosecond-class deterministic timing when coherent \IandQ processing must preserve stable range bins and integration phase.

\subsection{The Von Neumann Bottleneck at the PHY Boundary}
\textbf{Temporal execution model.}
Conventional CPUs execute a fetch-decode-execute stream over shared memory and interconnect resources. Latency is shaped by instruction flow, cache state, and branch behavior, which are optimized for average throughput rather than strict cycle-by-cycle guarantees.

\textbf{PHY mismatch points.}
\begin{itemize}
  \item Interrupt preemption introduces latency outliers.
  \item OS scheduler decisions create variable service intervals.
  \item Cache misses and bus arbitration cause nonuniform access latency.
  \item DMA completion timing can drift relative to sample-clock boundaries.
\end{itemize}
At the PHY layer, bounded worst-case latency is required; average latency is insufficient for deterministic sample-aligned operation.

\subsection{Spatial Computing on FPGA Fabric}
\textbf{Spatial definition.}
In FPGA gateware, functions are physically instantiated in fabric and evaluated each clock cycle. Throughput scales via parallel placement and pipelining rather than instruction retirement.

\textbf{Silicon hierarchy framing.}
\begin{center}
Logic Cell \(\rightarrow\) Slice \(\rightarrow\) CLB \(\rightarrow\) Column \(\rightarrow\) Die
\end{center}

\textbf{Determinism mechanism.}
Fixed routing, fixed pipeline depth, and fixed clock-domain constraints create predictable latency for \AXIStream dataflow through radar DSP stages.

\subsection{The Atomic Logic Cell: LUT and Flip-Flop}
\textbf{6-input LUT.}
A LUT6 is modeled as a \(2^6=64\)-entry truth-table memory. Each input vector selects a stored output value, implementing arbitrary combinational logic within timing and fanout limits.

\textbf{Flip-Flop as temporal barrier.}
The flip-flop samples LUT output on an active clock edge, creating a pipeline boundary that enables setup/hold analysis, retiming, and higher achievable \(F_{\max}\).

\begin{figure}[htbp]
  \centering
  \begin{tikzpicture}[node distance=2.4cm, >=Latex]
    \node[draw, minimum width=3.6cm, minimum height=2.0cm] (lut) {6-Input LUT\\(64-entry truth table)};
    \node[draw, right=of lut, minimum width=2.2cm, minimum height=2.0cm] (ff) {FF};

    \draw[->] ($(lut.west)+(-1.7,0.75)$) -- ($(lut.west)+(0,0.75)$) node[midway, above] {\(I_0\)};
    \draw[->] ($(lut.west)+(-1.7,0.45)$) -- ($(lut.west)+(0,0.45)$) node[midway, above] {\(I_1\)};
    \draw[->] ($(lut.west)+(-1.7,0.15)$) -- ($(lut.west)+(0,0.15)$) node[midway, above] {\(I_2\)};
    \draw[->] ($(lut.west)+(-1.7,-0.15)$) -- ($(lut.west)+(0,-0.15)$) node[midway, above] {\(I_3\)};
    \draw[->] ($(lut.west)+(-1.7,-0.45)$) -- ($(lut.west)+(0,-0.45)$) node[midway, above] {\(I_4\)};
    \draw[->] ($(lut.west)+(-1.7,-0.75)$) -- ($(lut.west)+(0,-0.75)$) node[midway, above] {\(I_5\)};

    \draw[->] (lut.east) -- node[above] {Combinational output} (ff.west);
    \draw[->] ($(ff.south)+(0,-0.9)$) -- ($(ff.south)+(0,0)$) node[midway, right] {clk \(\uparrow\)};
    \draw[->] (ff.east) -- ++(1.8,0) node[midway, above] {Registered output \(Q[n]\)};
    \node[anchor=west] at ($(ff.west)+(0.15,0.35)$) {\(D\)};
    \node[anchor=east] at ($(ff.east)+(-0.15,0.35)$) {\(Q\)};
  \end{tikzpicture}
  \caption{The LUT implements combinational mapping, while the flip-flop enforces cycle boundaries that make timing analyzable and deterministic.}
  \label{fig:lut_ff_timing_boundary}
\end{figure}

\subsection{Implications for Radar Front-End Precision}
\begin{itemize}
  \item Deterministic timing \(\Rightarrow\) stable range-bin assignment.
  \item Deterministic processing phase \(\Rightarrow\) coherent accumulation support.
  \item Bounded latency \(\Rightarrow\) repeatable trigger-to-capture behavior.
\end{itemize}
This section establishes the physical basis for subsequent gate-level resource and timing-closure discussions.

\bibliographystyle{ieeetr}
\bibliography{../shared/references}

\end{document}
